\documentclass[11pt]{article}

\usepackage[a4paper,margin=1in]{geometry}
\usepackage{amsmath,amssymb,amsthm}
\usepackage{algorithm2e}
\usepackage{enumitem}
\usepackage{hyperref}
\usepackage{xcolor}

\hypersetup{
  colorlinks=true,
  linkcolor=blue!45!black,
  citecolor=blue!45!black,
  urlcolor=blue!45!black
}

\newtheorem{definition}{Definition}
\newtheorem{remark}{Remark}
\numberwithin{equation}{section}

\newcommand{\Sym}{\operatorname{Sym}}
\newcommand{\GL}{\operatorname{GL}}
\newcommand{\rank}{\operatorname{rank}}
\newcommand{\Del}{\operatorname{Del}}
\newcommand{\cone}{\operatorname{cone}}
\newcommand{\linspan}{\operatorname{span}}
\newcommand{\ZZ}{\mathbb{Z}}
\newcommand{\RR}{\mathbb{R}}

\title{Algorithms for Computing Iso-Delaunay Domains}
\author{Mathieu Dutour Sikiri{\'c}}
\date{\today}

\begin{document}
\maketitle

\begin{abstract}
This note records an algorithm for enumerating iso-Delaunay domains, also
called \(T\)-secondary cones or \(L\)-type domains in a prescribed linear space
\(T\) of quadratic forms.  The first section gives the enumeration algorithm.
The second section explains a randomized search for an iso-Delaunay domain
whose extreme rays all define full-rank Gram matrices.
\end{abstract}

\section{Enumeration of iso-Delaunay domains}

\subsection{The input space}

Let \(T \subset \Sym_n(\RR)\) be the rational linear space of symmetric
matrices under consideration.  We use its saturated integral matrix lattice
\begin{equation}
  T_{\ZZ}=T\cap\Sym_n(\ZZ)
\end{equation}
and assume that it spans \(T\) over \(\RR\).  Fix a \(\ZZ\)-basis
\(B_1,\ldots,B_m\) of \(T_{\ZZ}\), and write
\(\ell_u(x)=x^{t}B_u x\).  Thus the integral coordinate vectors are exactly
the integral matrices in \(T\).  A point
\(q=(q_1,\ldots,q_m)\in\RR^m\) represents the Gram matrix
\begin{equation}
  Q(q)=\sum_{u=1}^{m} q_u B_u .
\end{equation}
The enumeration is performed inside the positive definite part of \(T\).  A
domain is represented by a Delaunay tessellation \(\Del(Q)\), an interior Gram
matrix \(Q\), and an invariant short-vector family used for equivalence tests.

\begin{definition}
An iso-Delaunay domain in \(T\) is the relatively open cone of all
positive-definite forms \(Q\in T\) that have a fixed Delaunay tessellation.
Its closure is a rational polyhedral cone in \(T\).
\end{definition}

In particular, if \(U\in\GL_n(\ZZ)\) preserves \(T\), then
\(Q\mapsto U^tQU\) preserves \(T_{\ZZ}\).  Its induced matrix in the above
\(T\)-coordinates therefore belongs to \(\GL_m(\ZZ)\).  This integral
coordinate structure is also used in the enumeration of \(T\)-perfect forms
\cite{Schuermann2009}.

The algorithm follows the \(T\)-secondary-cone theory of
\cite{DutourSikiriSchurmannVallentin2008}, itself a generalization of
Voronoi's reduction theory \cite{Voronoi1908}.  The finite object that is
enumerated is not the full infinite tessellation of \(\RR^n\), but a list of
orbits of Delaunay polytopes together with facet adjacencies and stabilizer
groups.

\subsection{Linear inequalities from one Delaunay wall}

The fundamental local computation is the construction of the wall inequality
associated with a Delaunay polytope \(P\) and a neighbouring Delaunay polytope
\(P'\).  Let \(v_1,\ldots,v_s\in\ZZ^n\) be the vertices of \(P\), and choose
an affine basis \(v_{i_0},\ldots,v_{i_n}\) of \(P\).  For a test vertex \(w\),
write \(w\) as an affine combination of this basis, that is,
\begin{equation}
  \sum_{k=0}^{n}\lambda_k = 1,\qquad
  w = \sum_{k=0}^{n}\lambda_k v_{i_k}.
\end{equation}
Here ``affine basis'' means only that the \(n+1\) selected vertices are
affinely independent, so the coefficients \(\lambda_k\) are real and uniquely
determined.  This differs from the terminology of Deza and Laurent
\cite{DezaLaurent1997}, where an affine basis of a lattice polytope is required
to give integral affine coefficients for every vertex of the polytope.  That
stronger condition cannot be imposed here: there exist Delaunay polytopes
having no affine basis in the integral sense
\cite{DutourSikiricGrishukhin2007}.
Then the Delaunay condition that \(w\) is outside or on the circumsphere of
the basis is the linear condition
\begin{equation}
  \sum_{u=1}^{m} q_u
  \left(
    \ell_u(w) - \sum_{k=0}^{n}\lambda_k \ell_u(v_{i_k})
  \right) \geq 0 .
\end{equation}
For an actual facet adjacency \(P,P'\), any vertex of \(P'\setminus P\) gives
the same inequality up to a scalar factor.  The scalar is canonicalized, and
all Delaunay-adjacency records producing that same wall are grouped together.

For a Delaunay tessellation \(D\), one loops over all orbits of Delaunay
polytopes and over all their facet adjacencies.  This gives
\begin{equation}
  a_1(q)\geq 0,\ldots,a_N(q)\geq 0,
\end{equation}
the full half-space description of the closure of the iso-Delaunay domain.
Redundant inequalities are then removed, leaving the irredundant wall
description used for adjacency.  This redundancy removal is carried out by
Clarkson's output-sensitive redundancy algorithm \cite{Clarkson1994}, using
the implementation available in cddlib \cite{FukudaCddlib}.

\begin{algorithm}\label{Algo_WallInequalities}
\KwData{A Delaunay tessellation \(D\) represented by orbit representatives of
Delaunay polytopes and their facet adjacencies; a basis
\(B_1,\ldots,B_m\) of \(T\)}
\KwResult{The irredundant wall inequalities of the iso-Delaunay cone of \(D\)}
${\mathcal A}\leftarrow \emptyset$.\\
\For{each orbit representative \(P\) of a Delaunay polytope of \(D\)}{
  choose an affine basis \(v_{i_0},\ldots,v_{i_n}\) of \(P\).\\
  \For{each facet-neighbour \(P'\) of \(P\)}{
    choose a vertex \(w\in P'\setminus P\).\\
    compute \(\lambda_0,\ldots,\lambda_n\) from (1.2).\\
    \(a(q)\leftarrow\) the linear form occurring in (1.3).\\
    replace \(a\) by its canonical scalar multiple.\\
    append to \(a\) the adjacency record \((P,P')\).\\
    ${\mathcal A}\leftarrow {\mathcal A}\cup\{a\}$, merging equal canonical
    inequalities and their adjacency records.\\
  }
}
remove redundant inequalities from \({\mathcal A}\) by Clarkson's algorithm.\\
\KwRet{\({\mathcal A}\)}
\caption{Wall inequalities of an iso-Delaunay cone}
\end{algorithm}

\subsection{Interior form of a domain}

Let
\begin{equation}
  C=\{q\in\RR^m : a_i(q)\geq 0,\ 1\leq i\leq N\}
\end{equation}
be the closure of an iso-Delaunay cone, assumed full-dimensional in its
ambient \(T\)-space.  The enumeration needs a representative point in the
relative interior of \(C\).  A generic linear program would find such a point,
but the result would depend on accidental choices in the linear program.  The
algorithm instead constructs a point that is canonical with respect to
permuting the inequalities and equivariant under linear changes of coordinates
in \(T\).

The first step fixes the scale of the cone by working on the affine section
\begin{equation}
  S=\{q\in\RR^m : a_i(q)\geq 1 \text{ for all } i\}.
\end{equation}
On this section minimize the symmetric objective
\begin{equation}
  \sum_{i=1}^{N} (a_i(q)-1),
\end{equation}
or, what is the same, minimize \(\sum_i a_i(q)\).  This objective is the sum of
all inequalities, so it is independent of their order.  If the optimum is a
single point, that point is returned.  It lies in the interior of \(C\), since
all inequalities in the scaled section are strict or equal to \(1>0\).

If the optimum is not a single point, the linear program has selected a
positive-dimensional optimal face
\begin{equation}
  F=\left\{q\in S : \sum_i(a_i(q)-1)=\mu\right\},
\end{equation}
where \(\mu\) is the optimal value.  One then works in the direction space
\begin{equation}
  L_F=\linspan\{q-q' : q,q'\in F\}.
\end{equation}
The inequalities that still cut the face nontrivially are projected to
coordinates on \(L_F\), and the same construction is applied recursively in
that lower-dimensional space.

The recursion strictly lowers the dimension whenever the first minimization is
degenerate.  At the last step either the symmetric linear program has a unique
solution or the remaining optimal face has no further nontrivial direction.
Pulling the recursive solution back through the successive affine projections
gives the final point \(q\in C\).

This construction uses only linear programming and linear algebra; no full
dual-description computation is needed.  It is geometrically canonical in the
following sense.  Permuting the inequalities does not change the point, and a
linear change of coordinates in \(T\) sends the selected point by the
contragredient transformation.  Rescaling the inequalities individually is not
an invariance of the construction, so a canonical normalization of the wall
inequalities is part of the preceding step.

\begin{algorithm}\label{Algo_GeomInterior}
\KwData{A full-dimensional cone \(C=\{q:a_i(q)\geq 0\}\), with the
inequalities \(a_i\) canonically scaled}
\KwResult{A geometrically canonical point in the relative interior of \(C\)}
\If{\(\dim C=0\)}{
  \KwRet{the unique point of \(C\)}
}
\(S\leftarrow \{q:a_i(q)\geq 1\mbox{ for all }i\}\).\\
Solve the linear program minimizing \(\sum_i(a_i(q)-1)\) on \(S\).\\
\(F\leftarrow\) the optimal face of this linear program.\\
\If{\(F\) is a single point \(q\)}{
  \KwRet{\(q\)}
}
\(L_F\leftarrow\linspan\{q-q':q,q'\in F\}\).\\
Choose an affine parametrization \(q=p+Uy\) of \(F\), with columns of \(U\)
spanning \(L_F\).\\
Project to the \(y\)-coordinates the inequalities that are not identically
constant on \(p+L_F\).\\
\(y_0\leftarrow\) result of Algorithm~\ref{Algo_GeomInterior} for the projected
system.\\
\KwRet{\(p+Uy_0\)}
\caption{Geometrically unique interior point}
\end{algorithm}

\subsection{Construction of an initial form}

It remains to construct a form lying in a full-dimensional iso-Delaunay
domain.  This is done iteratively.  Fix a positive-definite matrix
\(Q_{\mathrm{ref}}\in T\).  Choose integral coefficients \(c_u\), form
\begin{equation}
  Q=\sum_{u=1}^{m}c_uB_u,
\end{equation}
and add sufficiently many positive multiples of \(Q_{\mathrm{ref}}\) to make
\(Q\) positive definite.  If an arithmetic automorphism of \(Q\) does not fix
the \(T\)-space pointwise, discard \(Q\) and sample again.

The random generation starts with simple forms: the coefficients \(c_u\) are
sampled from a small interval around zero.  If repeated candidates fail the
tests below, this interval is progressively enlarged.  The precise rate of
growth is not part of the mathematical algorithm; its purpose is to explore
low-complexity forms first without restricting the search to a fixed finite
set.

For a form passing this symmetry test, enumerate its Delaunay tessellation.
This enumeration may be stopped as soon as a Delaunay polytope imposes a
nonzero linear equality on \(T\).  Such an equality means that \(Q\) lies on a
proper iso-Delaunay wall rather than in a full-dimensional domain.  If a
common form \(Q_{\mathrm{common}}\) is prescribed, stop also when a Delaunay
polytope of \(Q\) is not Delaunay for \(Q_{\mathrm{common}}\).

Once a tessellation \(D_0\) passes all tests, its wall inequalities are
computed and Algorithm~\ref{Algo_GeomInterior} supplies the distinguished
interior form \(Q_0\).  The iterative construction is summarized in
Algorithm~\ref{Algo_InitialIsoDelaunay}.

\begin{algorithm}\label{Algo_InitialIsoDelaunay}
\KwData{The saturated \(T\)-space with basis \(B_1,\ldots,B_m\), a
positive-definite form \(Q_{\mathrm{ref}}\in T\), and optionally
\(Q_{\mathrm{common}}\)}
\KwResult{A Delaunay tessellation \(D_0\) of a full-dimensional
iso-Delaunay domain and its distinguished interior form \(Q_0\)}
\While{no initial domain has been found}{
  Choose random integral coefficients \(c_1,\ldots,c_m\) according to the
  sampling policy described above.\\
  \(Q\leftarrow\sum_{u=1}^{m}c_uB_u\).\\
  \While{\(Q\) is not positive definite}{
    \(Q\leftarrow Q+Q_{\mathrm{ref}}\).
  }
  \If{some arithmetic automorphism of \(Q\) does not fix \(T\) pointwise}{
    discard \(Q\) and continue with another sampled form.
  }
  Enumerate the Delaunay tessellation \(D\) of \(Q\), stopping if a Delaunay
  polytope imposes a nonzero equality on \(T\).\\
  \If{\(Q_{\mathrm{common}}\) is prescribed}{
    stop the enumeration also if a Delaunay polytope is not Delaunay for
    \(Q_{\mathrm{common}}\).
  }
  \If{the enumeration finishes without stopping}{
    \({\mathcal A}_0\leftarrow\) wall inequalities of \(D\) by
    Algorithm~\ref{Algo_WallInequalities}.\\
    \(Q_0\leftarrow\) result of Algorithm~\ref{Algo_GeomInterior} applied to
    \({\mathcal A}_0\).\\
    \KwRet{\((D,Q_0)\)}
  }
}
\caption{Iterative construction of an initial iso-Delaunay domain}
\end{algorithm}

\subsection{Adjacent domains and flips}

The adjacency graph of iso-Delaunay domains is obtained by crossing
irreducible walls of the current cone.  For each orbit representative of an
irredundant inequality, the algorithm first computes a point in the relative
interior of the corresponding facet.  If this point does not represent a
positive-definite matrix, the wall belongs to the boundary of the positive
definite cone and no adjacent iso-Delaunay domain is produced.  Otherwise the
wall is crossed by a \(T\)-flip.

The flip uses the adjacency records stored with the wall inequality to identify
all Delaunay polytopes that meet at the wall.  These polytopes form
repartitioning polytopes.  The repartitioning construction lifts vertices by
the height \(x^tQx\) for an interior form \(Q\) of the old domain.  Its facets
split into three classes:
\begin{itemize}[leftmargin=2em]
  \item lower facets, which are the old Delaunay polytopes;
  \item higher facets, which become the new Delaunay polytopes after crossing
        the wall;
  \item barrel facets, which transmit adjacencies between neighbouring
        repartitioning polytopes.
\end{itemize}
After the higher facets have been identified modulo the local symmetry group,
one reconstructs the whole finite orbit description of the new Delaunay
tessellation and recomputes its facet adjacencies.

\begin{algorithm}\label{Algo_AdjacentIsoDelaunay}
\KwData{An iso-Delaunay domain \(C\), represented by its Delaunay tessellation
\(D\), an interior form \(Q\), and the wall inequalities of \(C\)}
\KwResult{The list of positive-definite iso-Delaunay domains adjacent to \(C\)}
\({\mathcal N}\leftarrow\emptyset\).\\
Compute the stabilizer of \(C\) in the arithmetic group preserving \(T\).\\
Compute orbit representatives of the irredundant walls of \(C\).\\
\For{each wall representative \(F\)}{
  compute a point \(q_F\) in the relative interior of \(F\).\\
  \If{\(Q(q_F)\) is not positive definite}{
    continue to the next wall.
  }
  collect all Delaunay adjacencies producing the wall \(F\).\\
  construct the associated repartitioning polytopes.\\
  replace lower facets by higher facets, keeping barrel facets for adjacency
  matching.\\
  reconstruct the finite orbit description of the flipped Delaunay
  tessellation \(D'\).\\
  compute an interior form \(Q'\) of the cone of \(D'\) by
  Algorithm~\ref{Algo_GeomInterior}.\\
  \({\mathcal N}\leftarrow{\mathcal N}\cup\{(D',Q')\}\).
}
\KwRet{\({\mathcal N}\)}
\caption{Adjacent iso-Delaunay domains}
\end{algorithm}

\subsection{Equivalence and graph traversal}

Enumeration is a graph traversal in which each visited vertex is an
iso-Delaunay domain and each edge is a positive-definite wall crossing.  The
data required by this traversal are:
\begin{itemize}[leftmargin=2em]
  \item construction of one initial iso-Delaunay domain;
  \item computation of all adjacent domains by the wall-and-flip procedure
        above;
  \item computation of a fast invariant, for instance from an invariant
        short-vector family of the interior Gram matrix;
  \item an exact equivalence test between two domains, obtained by finding a
        \(T\)-space-preserving integral change of basis sending one interior
        Gram matrix and short-vector family to the other.  Because the matrix
        lattice \(T_{\ZZ}\) is saturated, the induced transformation matrix
        in the chosen \(T\)-coordinates is integral.
\end{itemize}
The output of the traversal is therefore one representative for each
iso-Delaunay domain modulo the arithmetic group preserving the \(T\)-space.

\subsection{Enumeration algorithm}

The top-dimensional enumeration is obtained by putting the preceding
subroutines together.  Its structure is the usual adjacency enumeration:
maintain a queue of domains still to be processed, compute all neighbours of
one domain, and insert only those neighbours that are not equivalent to
previously found domains.

\begin{algorithm}\label{Algo_IsoDelaunayEnumeration}
\KwData{A \(T\)-space with basis \(B_1,\ldots,B_m\), an arithmetic group
preserving \(T\), and optionally a common Gram matrix \(Q_{\mathrm{common}}\)}
\KwResult{A set \({\mathcal R}\) of inequivalent full-dimensional
iso-Delaunay domains in \(T\)}
\((D_0,Q_0)\leftarrow\) result of
Algorithm~\ref{Algo_InitialIsoDelaunay}.\\
\({\mathcal T}\leftarrow\{(D_0,Q_0)\}\), \({\mathcal R}\leftarrow\emptyset\).\\
\While{\({\mathcal T}\neq\emptyset\)}{
  choose \(C=(D,Q)\in{\mathcal T}\).\\
  \({\mathcal T}\leftarrow{\mathcal T}\setminus\{C\}\).\\
  \({\mathcal R}\leftarrow{\mathcal R}\cup\{C\}\).\\
  \({\mathcal N}\leftarrow\) adjacent domains of \(C\) by
  Algorithm~\ref{Algo_AdjacentIsoDelaunay}.\\
  \For{\(C'\in{\mathcal N}\)}{
    \If{\(C'\) is not equivalent to any domain in \({\mathcal R}\cup{\mathcal T}\)}{
      \({\mathcal T}\leftarrow{\mathcal T}\cup\{C'\}\).
    }
  }
}
\KwRet{\({\mathcal R}\)}
\caption{Enumeration of full-dimensional iso-Delaunay domains}
\end{algorithm}

\section{Finding iso-Delaunay domains with all rays of full rank}

The following search is more targeted than full enumeration.  Its purpose is to
find, in a fixed dimension \(n\), a top-dimensional iso-Delaunay domain whose
every extreme ray gives a full-rank Gram matrix.  The search is not
dimension-specific.

\subsection{The objective function}

Let \(C\subset T\) be the closure of an iso-Delaunay domain with irredundant
inequality description.  A dual-description computation returns its extreme
rays
\begin{equation}
  r^{(1)},\ldots,r^{(s)}\in \RR^m .
\end{equation}
Each ray is interpreted as a symmetric matrix
\begin{equation}
  R_i=\sum_{u=1}^{m} r^{(i)}_u B_u .
\end{equation}
The desired property is
\begin{equation}
  \rank(R_i)=n \qquad \text{for every extreme ray } i .
\end{equation}
The search uses the integer objective
\begin{equation}
  \Phi(C)=\sum_{i=1}^{s} (n-\rank(R_i)).
\end{equation}
Thus \(\Phi(C)=0\) exactly when all extreme rays have full rank.  For each
visited domain, the inequalities are recomputed, redundancy is removed by
Clarkson's algorithm, the extreme rays are computed, and the rank defects are
summed.  This makes the objective independent of any redundant wall
description used during adjacency generation.

\begin{algorithm}\label{Algo_RankDefect}
\KwData{An iso-Delaunay domain \(C\) with Delaunay tessellation \(D\)}
\KwResult{The rank-defect score \(\Phi(C)\)}
Compute the wall inequalities of \(D\) by Algorithm~\ref{Algo_WallInequalities}.\\
Remove redundant inequalities by Clarkson's algorithm.\\
Compute the extreme rays \(r^{(1)},\ldots,r^{(s)}\) of the resulting cone.\\
\(\Phi\leftarrow 0\).\\
\For{\(i=1,\ldots,s\)}{
  \(R_i\leftarrow\sum_{u=1}^m r^{(i)}_u B_u\).\\
  \(\Phi\leftarrow \Phi+(n-\rank(R_i))\).
}
\KwRet{\(\Phi\)}
\caption{Rank-defect objective}
\end{algorithm}

\subsection{Greedy descent with random escapes}

The search starts from a generic iso-Delaunay domain obtained by
Algorithm~\ref{Algo_InitialIsoDelaunay}.  The strategy is the same one used
in Frank Vallentin's PhD thesis \cite{VallentinPhD}: move greedily in the
graph of adjacent \(L\)-types when the objective improves, and use random walks
to leave local minima.  In that thesis this strategy was used to find record
lattice coverings in dimension \(6\).

\begin{algorithm}\label{Algo_FullRankSearch}
\KwData{An \(n\)-dimensional \(T\)-space, a random-walk length
\(N_{\rm walk}\), and an iteration bound \(N_{\rm iter}\)}
\KwResult{A domain \(C\) with \(\Phi(C)=0\), if one is found within the search
budget}
\(C\leftarrow\) result of Algorithm~\ref{Algo_InitialIsoDelaunay}.\\
\(\phi\leftarrow\Phi(C)\), computed by Algorithm~\ref{Algo_RankDefect}.\\
\For{\(t=1,\ldots,N_{\rm iter}\)}{
  \If{\(\phi=0\)}{
    \KwRet{\(C\)}
  }
  \({\mathcal N}\leftarrow\) adjacent domains of \(C\) by
  Algorithm~\ref{Algo_AdjacentIsoDelaunay}.\\
  compute \(\Phi(C')\) for all \(C'\in{\mathcal N}\).\\
  \(\phi_{\min}\leftarrow \min\{\Phi(C'):C'\in{\mathcal N}\}\).\\
  \eIf{\(\phi_{\min}<\phi\)}{
    choose uniformly a domain \(C'\in{\mathcal N}\) with
    \(\Phi(C')=\phi_{\min}\).\\
    \(C\leftarrow C'\), \(\phi\leftarrow\phi_{\min}\).\\
  }{
    perform \(N_{\rm walk}\) random adjacency steps from \(C\).\\
    replace \(C\) by the endpoint and recompute \(\phi=\Phi(C)\).
  }
}
\KwRet{\bf fail}
\caption{Search for an iso-Delaunay domain with full-rank rays}
\end{algorithm}

The output condition is exactly \(\Phi(C)=0\).  In this case every extreme ray
of the iso-Delaunay domain has full rank.

\subsection{Search parameters}

The search is controlled by the following parameters:
\begin{itemize}[leftmargin=2em]
  \item the length of the random escape walk after a local minimum;
  \item the cap on the number of outer iterations.
\end{itemize}

The iteration cap is part of the mathematical design of the search, not only a
safety guard.  In some dimensions there may be no domain with \(\Phi=0\), so an
unconditional search for a zero would not terminate.  The search exploits the
adjacency graph of iso-Delaunay domains until such a domain is found or the
chosen computational budget is exhausted.

\nocite{DutourSikiricPolyhedralCommon}
\bibliographystyle{plain}
\bibliography{references}

\end{document}
